\begin{proof}

\begin{enumerate}[label={$\langle1\rangle$\arabic*.},leftmargin=*]
\item We first identify exactly the analytic comparison used in the
argument. Put
\[
\mathscr K^\bullet=
 \Omega_U^\bullet(\log D)\big/
 \bigl(h^*\Omega_\Delta^1(\log0)\wedge
       \Omega_U^{\bullet-1}(\log D)\bigr),\qquad
\mathscr H^2=\mathbb R^2h_*\mathscr K^\bullet.
\]
In \cite[\S2.1, p.~234; Theorem~2.18, pp.~237--238; \S2.19 and Proposition~2.20, pp.~238--239; Theorem~2.21, p.~239]{Steenbrink1976}, the assumptions for these assertions
are precisely a proper holomorphic map from a connected complex
manifold to a disc, smooth off zero, with normal-crossings support.
They precede the additional Hodge-theoretic assumptions of the
later sections. They give a locally free bundle $\mathscr H^2$,
its actual Gauss--Manin connection
\[
\nabla:\mathscr H^2\longrightarrow
\Omega_\Delta^1(\log0)\otimes\mathscr H^2,
\]
and the comparison
$\mathscr H^2|_{\Delta^*}\simeq
R^2h_*\C\otimes\OO_{\Delta^*}$.
The residue eigenvalues belong to $[0,1)$ and the extended
monodromy at zero is
$\exp(-2\pi i\operatorname{Res}_0\nabla)$.
Its canonical base-change map identifies the fibre at zero with
$\mathbb H^2(D,\mathscr K^\bullet/t\mathscr K^\bullet)$.
The local calculation used below is the same monomial Koszul
calculation as \cite[Proposition~1.13, pp.~233--234, and proof of Proposition~2.20, pp.~238--239]{Steenbrink1976}; we give it explicitly for the
subcomplex needed here. No local invariant-cycle surjectivity
or polarization is invoked.

\item Suppose that the monodromy on $H^2$ were the identity. Its
residue has all eigenvalues zero, and its exponential is the
identity. A nilpotent matrix whose exponential is the identity
is zero. Hence the residue vanishes. In a holomorphic frame
$\nabla=d+A(t)\,dt/t$ with $A(0)=0$, so this is a holomorphic
connection across zero. It has a holomorphic horizontal frame
on the disc. Consequently the coordinates of every holomorphic
section of $\mathscr H^2$ in a flat frame, that is its periods,
are holomorphic and bounded near zero.

\item The determinant of the locally free relative logarithmic complex is
\[
\mathscr K^2\simeq
 \omega_U(D)\otimes h^*\omega_\Delta(0)^{-1}
 \simeq\omega_{U/\Delta}(D-F).
\]
The direct image $h_*\mathscr K^2$ is a torsion-free
coherent sheaf of rank one, because a smooth torus fibre has
one-dimensional space of holomorphic two-forms. Shrink the disc
and choose a generator $\Omega$ of this line bundle.
Thus $\Omega$ determines a holomorphic section $[\Omega]$ of
$\mathscr H^2$ by the inclusion of the last term of the
relative de Rham complex. Its smooth-fibre comparison is the
ordinary class $[\Omega_t]\in H^2(U_t;\C)$.

\item First suppose that $\Omega$ restricts nontrivially to a normalized
component $T$ of $D$. There is a canonical isomorphism
\[
 \mathscr K^1|_T\simeq\Omega_T^1(\log D_T),\qquad
 \mathscr K^2|_T\simeq\omega_T(D_T),
\]
where $D_T$ is the reduced conductor. To specify the map, on the
branch $z_i=0$ subtract
$\operatorname{res}_i(\eta)(dt/t)/m_i$ from an absolute logarithmic
one-form $\eta$ and restrict its tangential part. The normal
residue of $dt/t$ is $m_i$, so this kills exactly the subbundle
divided out in $\mathscr K^1$ and is independent of the lift.
The exterior square gives the second isomorphism. If every
residue of $\Omega|_T$ vanished, its local logarithmic expression
would be holomorphic on $T$. This contradicts $p_g(T)=0$.
Thus there is a smooth point of a double curve, away from triple
points, where its residue is nonzero.

\item In a small chart about that point, write $h=x^ay^b=t$ and use $z$
along the double curve. The relative two-form is
\[
 \Omega_t=g(x,y,z)\,\frac{dx}{x}\wedge dz,
 \qquad g(0,0,z_0)\ne0.
\]
Choose a fixed small $z$-disc and $\epsilon>0$ such that
$|g(x,y,z)|\ge c>0$ whenever $|x|,|y|<\epsilon$ there.
The map to the annulus
$|t|^{1/a}\epsilon^{-b/a}<|x|<\epsilon$
has $b$ sheets. Integration on its inverse image and the fixed
$z$-disc gives
\[
 \int_{U_t}\Omega_t\wedge\overline{\Omega_t}
 \ \ge\ c'\log\frac{\epsilon^{a+b}}{|t|}
 \ \longrightarrow\ +\infty.
\]
Here the integrand is a nonnegative real volume form on the
complex surface; its local expression is a positive constant
times $|g|^2|x|^{-2}\,dA(x)dA(z)$.

\item On the other hand, $\Omega_t$ is closed, and this integral is
the fixed cup-product pairing of $[\Omega_t]$ with its complex
conjugate. All its periods are bounded by the holomorphic
connection argument, so the pairing is bounded. This is a
contradiction.

\item It remains to treat the case where $\Omega$ vanishes on every
branch of $D$. Put
\[
 \mathscr J=\OO_U(-D),\qquad
 \mathscr B^\bullet=
 \mathscr J\mathscr K^\bullet/t\mathscr K^\bullet.
\]
Both $\mathscr J\Omega_U^\bullet(\log D)$ and
$t\Omega_U^\bullet(\log D)$ are subcomplexes for the absolute
exterior derivative: locally the logarithmic derivatives of
$z_1\cdots z_r$ and $t=\prod z_i^{m_i}$ are holomorphic
linear combinations of the $d\log z_i$. Hence the locally split
absolute/relative logarithmic exact sequence remains an exact
sequence of complexes after tensoring by $\mathscr J/(t)$:
\[
0\longrightarrow\mathscr B^\bullet[-1]
 \xrightarrow{\ (dt/t)\wedge\ }
 (\mathscr J/(t))\Omega_U^\bullet(\log D)
 \longrightarrow\mathscr B^\bullet\longrightarrow0.
\tag{D.18}\hypertarget{tag-0195}{}\label{cdp:eq:D.18}
\]
Signs in the shifted differential are chosen so that the first
arrow is a morphism of complexes. Its connecting morphism
$\rho_{\mathscr B}:\mathscr B^\bullet\to\mathscr B^\bullet$
in the derived category defines degree-preserving endomorphisms
on both cohomology sheaves and hypercohomology. The inclusion
$\mathscr J/(t)\hookrightarrow\OO_U/(t)$ gives a commutative
diagram from \eqref{cdp:eq:D.18} to the analogous exact sequence for
$\mathscr K^\bullet/t\mathscr K^\bullet$. The induced
endomorphism for the latter is precisely the Gauss--Manin
residue. Thus the map on hypercohomology intertwines these
endomorphisms.

\item We compute the eigenvalues of $\rho_{\mathscr B}$ without
assuming reduced multiplicities. At a point of $r$ branches,
write
\[
 t=\prod_{i=1}^r z_i^{m_i},\qquad
 e=\gcd(m_1,\ldots,m_r),\qquad
 y=\prod_{i=1}^r z_i^{m_i/e},\qquad t=y^e.
\]
The transverse holomorphic Poincar\'e lemma reduces the relative
complex to the Koszul complex with commuting operators
\[
 D_i=z_i\partial_{z_i}
           -\frac{m_i}{m_1}z_1\partial_{z_1},
 \qquad 2\le i\le r,
\]
on the convergent-power-series module
$(z_1\cdots z_r)/(z_1^{m_1}\cdots z_r^{m_r})$.
The monomials diagonalize these operators. A monomial with some
nonzero eigenvalue has a contracting Koszul homotopy, obtained
by division by that eigenvalue. These divisions preserve
convergence: every nonzero eigenvalue has absolute value at
least $1/m_1$. The surviving monomials have exponent vector
$a(m_1/e,\ldots,m_r/e)$. Membership in $\mathscr J$ requires
$a\ge1$, and quotienting by $t$ requires $a<e$. Consequently
each cohomology stalk is spanned by
\[
 y^a\,d\log z_{i_1}\wedge\cdots\wedge d\log z_{i_q},
 \qquad 1\le a<e,
\]
subject to the relative relation $\sum m_i d\log z_i=0$.
Since $dy^a=(a/e)y^a\,dt/t$, the connecting endomorphism
acts on these representatives by $a/e$. In particular zero is
not an eigenvalue; when $e=1$ the complex is acyclic.

\item The hypercohomology spectral sequence
\[
 H^p\bigl(D,\mathcal H^q(\mathscr B^\bullet)\bigr)
 \ \Longrightarrow\ \mathbb H^{p+q}(D,\mathscr B^\bullet)
\]
is equivariant for this derived endomorphism. There are only
finitely many branch multiplicities, so all eigenvalues on its
terms and abutment lie in a finite subset of $(0,1)$. For
example the product of the corresponding linear factors
annihilates every cohomology sheaf, and a power annihilates the
abutment. Hence $\rho_{\mathscr B}$ is invertible on
$\mathbb H^2(D,\mathscr B^\bullet)$. The residue on
$\mathscr H^2|_0$ is zero. Equivariance therefore proves that
\[
 \mathbb H^2(D,\mathscr B^\bullet)\longrightarrow
 \mathbb H^2(D,\mathscr K^\bullet/t\mathscr K^\bullet)
       =\mathscr H^2|_0
 \quad\hbox{is the zero map}.
\tag{D.19}\hypertarget{tag-0196}{}\label{cdp:eq:D.19}
\]

\item The central class of $\Omega$ comes from the left side of
\eqref{cdp:eq:D.19}, because $\Omega\in\mathscr J\mathscr K^2$.
Thus $[\Omega]$ vanishes at zero in the locally free bundle
$\mathscr H^2$, and $[\Omega]/t$ is a holomorphic section of
that bundle. Its periods are bounded. On a smooth fibre this
section is the de Rham class of $\Omega_t/t$.

\item Primitivity of the generator $\Omega$ implies that its order
$b_i$ along at least one $D_i$ is less than $m_i$; otherwise
$\Omega/t$ would be a holomorphic section of $\mathscr K^2$
over $U$, contradicting its choice as a generator of
$h_*\mathscr K^2$. Choose a smooth point of that component
away from all other branches and all further zeros. There
are coordinates with $t=z^{m_i}$ in which
\[
 \Omega=z^{b_i}a(z,\xi,\eta)\,d\xi\wedge d\eta,
 \qquad a(0,0,0)\ne0.
\]
On one sheet above a fixed small $(\xi,\eta)$-polydisc, the
positive integral of $\Omega_t/t$ against its conjugate has
the lower bound
\[
 \int_{U_t}\frac{\Omega_t}{t}\wedge
                    \overline{\frac{\Omega_t}{t}}
 \ \ge c''|t|^{2(b_i/m_i-1)}\longrightarrow+\infty.
\]
Its expression as the fixed cup-product pairing of bounded
periods again gives a contradiction. This proves nontriviality
of the $H^2$ monodromy in the remaining case.

\item Finally the natural
cup-product isomorphism
$\bigwedge^2H^1(U_t;\Q)\simeq H^2(U_t;\Q)$ for a two-torus
is monodromy-equivariant, proving the final assertion.

\end{enumerate}
\end{proof}